% Generated draft: requires mathematical and source review.
\documentclass[11pt]{article}
\usepackage[margin=25mm]{geometry}
\usepackage{fontspec}
\setmainfont{Latin Modern Roman}
\usepackage{amsmath,amssymb,mathtools,unicode-math}
\setmathfont{Latin Modern Math}
\usepackage{enumitem,needspace,xurl,hyperref}
\setlength{\parindent}{0pt}
\setlength{\parskip}{5pt}
\setlength{\emergencystretch}{3em}
\newcommand{\N}{\mathbb{N}}
\newcommand{\Z}{\mathbb{Z}}
\newcommand{\Q}{\mathbb{Q}}
\newcommand{\R}{\mathbb{R}}
\newcommand{\C}{\mathbb{C}}
\newcommand{\F}{\mathbb{F}}
\newcommand{\A}{\mathbb{A}}
\newcommand{\PP}{\mathbb{P}}
\newcommand{\E}{\mathbb{E}}
\newcommand{\eps}{\varepsilon}
\newcommand{\dd}{\,\mathrm{d}}
\begin{document}
\section*{UnsolvedMath ID: 30000762}
\textbf{Source problem code:} catalogue entry
\subsection{Definitions and mathematical statement}
Let \((\mathbb H^2,g)\) be the hyperbolic plane. An entire minimal graph in \(\mathbb H^2\times\R\) is the graph \(\Gamma_u=\{(x,u(x)):x\in\mathbb H^2\}\) of a smooth function \(u:\mathbb H^2\to\R\) satisfying
\[
 \operatorname{div}_g\left(\frac{\nabla_g u}{\sqrt{1+|\nabla_g u|_g^2}}\right)=0.
\]
Entire means that the domain is all of \(\mathbb H^2\), not that \(u\) is holomorphic. The general request is to describe asymptotic values of such \(u\) at the ideal boundary \(\partial_\infty\mathbb H^2\) and to establish appropriate Fatou-type boundary theorems. The supplied record does not specify whether the requested asymptotic values concern arbitrary curves, radial approaches, or more restricted approach regions, or whether infinite values are included. For example, a curve-asymptotic value at \(\zeta\) would be a limit of \(u(\gamma(t))\) along a curve \(\gamma(t)\to\zeta\); different curves need not be assumed to give the same limit.

The explicit bounded case can be formulated using a base point \(o\in\mathbb H^2\). Let \(\gamma_\theta:[0,\infty)\to\mathbb H^2\) be the unit-speed geodesic ray from \(o\) in direction \(\theta\in S_o^1\), and use angular Lebesgue measure on \(S_o^1\), equivalently the corresponding visual measure on the ideal boundary. The reported conjecture asks whether
\[
 \left[\sup_{x\in\mathbb H^2}|u(x)|<\infty\right]
 \quad\Longrightarrow\quad
 \lim_{r\to\infty}u(\gamma_\theta(r))\text{ exists in }\R
 \text{ for almost every }\theta\in S_o^1.
\]
In a disk model centered at \(o\), this is the existence of \(\lim_{\rho\uparrow1}u(\rho e^{i\theta})\) for almost every boundary angle. It asserts radial limits, not automatically limits along every approach to a boundary point.

The supplied assessment dated 2026-08-21 reports that Collin--Rosenberg proved almost-everywhere radial limits for broad classes of minimal graphs on the disk, including a hyperbolic setting, while the full asymptotic-value classification and all Fatou-type phenomena remain unresolved. The precise hypotheses and conclusions of that reported progress are not reproduced here.
\subsection{Short English statement}
The broad question asks for boundary behavior of globally defined minimal graph functions over the hyperbolic plane. Its concrete bounded case asks for a finite limit along almost every geodesic ray. Supplied later work is reported as partial progress; radial, unrestricted, and other boundary approaches must be distinguished.
\subsection{Sources}
\begin{itemize}
\item Local UnsolvedMath registry, permanent ID \(30000762\), exact source code \texttt{catalogue entry}, proposed year 2007; no individual proposer is supplied.
\item \emph{Progress in Surface Theory (2008)}, Oberwolfach Reports, DOI 10.4171/owr/2007/24, pp. 1299--1376. \url{https://doi.org/10.4171/owr/2007/24}.
\item Supplied literature reference: Collin--Rosenberg, \emph{Annales de l'Institut Fourier}, \url{https://doi.org/10.5802/aif.2610}.
\item The \texttt{corrected\_verified} note identifies explicit asymptotic-value and radial-limit questions. The registry's partially solved label and assessment checked 2026-08-21 are dated supplied triage requiring independent verification of status and cited theorems.
\end{itemize}
\end{document}
